\documentclass[10pt,a4paper]{amsart}
\usepackage[margin=19mm]{geometry}
\usepackage{amsmath,amssymb,mathtools}
\usepackage[scr=rsfs,cal=euler]{mathalfa}
\usepackage{xfrac}
\usepackage[unicode,hidelinks,bookmarks=false]{hyperref}
\newcommand{\ie}{i.e.\ }
\newcommand{\cf}{cf.\ }
\newcommand{\resp}{respectively\ }
\newcommand{\Subsection}{\subsection}
\providecommand{\nobreakdash}{\nobreak\hbox{-}}
\setlength{\emergencystretch}{2em}
\multlinegap0pt
\usepackage{mathtools}
\usepackage{float}
\usepackage{amssymb}
\usepackage{enumerate}
\usepackage[shortlabels]{enumitem}
\usepackage{bm}
\usepackage{tikz}
\newtheorem{proposition}{Proposition}
\title{Localization and the Floer homology of strongly invertible knots}
\author{Aakash Parikh}
\date{Source: arXiv:2408.13892v2 (CC BY 4.0)}
\begin{document}
\maketitle

\noindent\emph{Source and licence.} Localization and the Floer homology of strongly invertible knots. Version arXiv:2408.13892v2 (CC BY 4.0), distributed by arXiv under the Creative Commons Attribution 4.0 International licence, http://creativecommons.org/licenses/by/4.0/, as recorded in the arXiv licence metadata for this exact version and shown on the abstract page https://arxiv.org/abs/2408.13892v2. Full licence notice: \texttt{licenses/arxiv-2408.13892-CC-BY-4.0.txt}.\par\smallskip

\noindent\textit{Editorial scope.} Counted unit: the article's proposition coM'/X'freeabelianfiniterank (source0.tex lines 1032--1034). The article's complete original proof of it is delivered with it (source0.tex lines 1107--1115, one \texttt{proof} environment of 9 lines, which the article places away from the statement). No mathematics is written, repaired or re-derived: the statement and the proof are the article's own lines, in the article's own notation, and no retained line is substituted or re-flowed.  No further source-local context is needed: every cross-reference of every retained line resolves inside the two delivered spans.  Nothing was omitted from any retained span, so the package carries no omission record.  No retained line contains a citation, so the wrapper carries no bibliography and no bibliography entry.  No retained line contains a figure environment, an image-inclusion command or drawing code, so no image is delivered and the package declares no asset dependency.  Editorial formatting only, in addition: an AMS \texttt{amsart} preamble that restates the article's own declarations and macros used by the retained lines, namely the environments \texttt{proposition} on separate counters, together with the standard packages those lines need.  The retained environments print in this document as Proposition 1 in order of appearance, whereas the article prints the same items with the numbering of its own full document.  The complete native source of the article as distributed in the arXiv e-print for this exact version is bundled unchanged as \texttt{sources/164253/source0.tex}.
\begin{proposition}\label{coM'/X'freeabelianfiniterank}
    The reduced cohomology $\widetilde{H}^*(M'\times[0,1]/X')$ is a free abelian group of finite rank.
\end{proposition}

\begin{proof}[Proof of Proposition \ref{coM'/X'freeabelianfiniterank}.]
In the following we use $Q$ instead of $q$ because the lowercase letter already denotes the quotient map $S^3\to S^3/\tau$. Surjectivity of $I^*$ for $m\ge 1$ tells us that the cohomology long exact sequence of the pair $(M'\times[0,1],X')$
\[\hdots \to\widetilde{H}^m(M'\times[0,1]/X')\xrightarrow{Q^*}\widetilde{H}^m(M'\times[0,1])\xrightarrow{I^*}\widetilde{H}^m(X')\xrightarrow{\partial}\widetilde{H}^{m+1}((M'\times[0,1]/X')\to\hdots\]
breaks up into short exact sequences
\[
0\to\widetilde{H}^m(M'\times[0,1]/X')\xhookrightarrow{}\widetilde{H}^m(M'\times[0,1])\twoheadrightarrow\widetilde{H}^m(X')\to 0
\]
for $m\ge 2$. So, for $m\ge 2$, $\widetilde{H}^m(M'\times[0,1]/X')$ injects into the free abelian group $\widetilde{H}^m(M'\times[0,1])$, and hence it is free abelian itself. All that is left is to analyze $\widetilde{H}^1(M'\times[0,1]/X')$, but the first cohomology of a space with finitely generated homology is always free.
\end{proof}

\end{document}
